% ====================================================================
% §2 통계역학 기초 (Part 0) — 파트 b: 평균장·전기화학 결선·거시 (2.4–2.6)
%   (앞 파일 ch1_sec02a_part0 의 \section 아래로 이어지는 subsection)
% ====================================================================
\subsection{평균장 상호작용 --- 정규용액 $\mu(\theta)$·$g(\xi)$ 와 이중웰(double-well) 문턱}\label{sec:sm-mf}
\textbf{(a) 출발 --- 이웃 상호작용의 평균장 근사.} 실제 격자에서 자리들은 독립이 아니다 --- 점유된 이웃은 다음
Li 의 삽입 에너지를 바꾼다(정전 반발·탄성 변형·전자 구조). 점유 이웃 \emph{쌍}당 에너지를 $u$, 배위수를 $z$
라 하면 정확한 상호작용 에너지는 배치마다 다르지만, 각 자리의 이웃 점유를 평균 $\theta$ 로 대체하는
\emph{평균장}(mean-field, Bragg--Williams) 근사에서
\begin{equation}
E_\mathrm{int}\;\simeq\;\frac{Mz}{2}\,u\,\theta^2
\label{eq:sm-mf}
\end{equation}
($M\theta$ 개 점유 자리 각각의 $z\theta$ 개 점유 이웃, $\tfrac12$ 은 쌍 이중계수 방지).

\textbf{(b) 연산 --- 항등 분해와 $\Omega$ 의 정의.} 자리 1몰당 상호작용 몫은 $\tfrac{z}{2}N_Au\,\theta^2$ 이고,
항등식 $\theta^2=\theta-\theta(1-\theta)$ 로 분해하면
\begin{equation}
\frac{z}{2}N_Au\,\theta^2
=\underbrace{\frac{z}{2}N_Au\,\theta}_{\text{선형 --- }\mu^0\text{ 에 흡수}}
+\;\Omega\,\theta(1-\theta),
\qquad
\Omega\;\equiv\;-\frac{z}{2}N_A\,u,
\label{eq:sm-omega}
\end{equation}
곧 점유 이웃끼리 \emph{인력}($u<0$)이면 $\Omega>0$(같은 종이 뭉치려는 상분리 경향), 반발($u>0$)이면
$\Omega<0$(교대 정렬 경향)이다. 이 한 모수로 혼합 자유에너지를 적는 모형이
\emph{정규용액}이고, 본론 표~\ref{tab:staging} 의 $\Omega_j$ 와 Part II 의
$\Omega_j^\mathrm{cat}$(표 밖 --- 피팅 배정 전제, \S\ref{sec:lco-hys}) [J/mol]이 정확히 이 슬롯이다.

\textbf{(c) 중간식 --- $g(\theta)$ 와 $\mu(\theta)$.} 자리 1몰당 자유에너지에 \eqref{eq:sm-Smix} 의 섞임 몫과
\eqref{eq:sm-omega} 의 상호작용 몫을 더하면(이하 한 전이의 $\Omega\equiv\Omega_j$)
\begin{equation}
g(\theta)=\mu^0\theta+RT\big[\theta\ln\theta+(1-\theta)\ln(1-\theta)\big]+\Omega_j\,\theta(1-\theta),
\label{eq:sm-gtheta}
\end{equation}
$\theta$ 로 미분하면(로그 몫 $RT\ln[\theta/(1-\theta)]$ --- 식~\eqref{eq:sm-mucount} 와 같은 계산; 상호작용 몫
$\Omega_j(1-2\theta)$, 상수 $\pm1$ 상쇄)
\begin{equation}
\mu_\mathrm{Li}(\theta)=\mu^0+RT\ln\frac{\theta}{1-\theta}+\Omega_j\,(1-2\theta).
\label{eq:mu}
\end{equation}
(그림~\ref{fig:sm-mu} 가 이 함수의 세 얼굴이다.) 진행률 $\xi=1-\theta$ 로 옮기면 섞임 몫과
$\theta(1-\theta)=\xi(1-\xi)$ 몫이 모두 $\xi\leftrightarrow1-\xi$ 자리바꿈에 대칭(우함수)이라 조성 몫의 꼴이
불변이고, 상수·직선 몫은 공통 접선·곡률 판정에 기여하지 않으므로 $g_j^0$(상수와 $\xi$-직선 몫의 묶음 ---
미분에는 상수 기여만 남긴다)로 모아 적는다:
\begin{equation}
g_j(\xi)=g_j^0+RT\big[\xi\ln\xi+(1-\xi)\ln(1-\xi)\big]+\Omega_j\,\xi(1-\xi),
\label{eq:gxi}
\end{equation}
--- 이 우함수 대칭이 \S\ref{sec:hys} 의 1계 미분 논거($f'(1-\xi)=-f'(\xi)$)의 뿌리다.

\textbf{(d) 박스 --- 볼록성 상실의 문턱.} 식~\eqref{eq:gxi} 를 두 번 미분하면
$g_j''(\xi)=RT/[\xi(1-\xi)]-2\Omega_j$ 이고, 첫 항의 최솟값이 $\xi=\tfrac12$ 에서 $4RT$ 이므로
\begin{equation}
\boxed{\;
\begin{aligned}
&g_j''(\xi)\ge0\ \ \forall\xi\ \Longleftrightarrow\ \Omega_j\le2RT\ \ (\text{단상·연속 고용체; 등호는 }\Omega_j{=}2RT,\ \xi{=}\tfrac12\ \text{한 점}),\\
&\Omega_j>2RT\ \Longleftrightarrow\ \text{이중웰}\ \ \Big(T_{c,j}=\frac{\Omega_j}{2R}\Big)\;.
\end{aligned}}
\label{eq:sm-thresh}
\end{equation}
$\Omega_j>2RT$ 면 가운데($g''<0$)가 언덕이 되어 두 골짜기가 생기고(그림~\ref{fig:sm-gxi}), 그 변곡점
(spinodal)의 닫힌 근과 거기서 자라는 히스테리시스 gap 은 \S\ref{sec:hys} 가 닫는다 --- Part 0 은 문턱까지만
놓는다.
\begin{verifybox}
극한 검산 --- $\Omega_j\to0$: $g''\ge4RT>0$ 로 항상 한 웰(\S\ref{sec:sm-lattice} 로 환원);
$\Omega_j\to2RT^-$: $g''(\tfrac12)\to0^+$ 로 바닥이 평탄해지되 아직 한 웰; $\Omega_j\to2RT^+$: 웰이 둘로
갈라지기 시작하고 $T_{c,j}=\Omega_j/2R$ 아래에서만 상분리가 존재한다(안정하다). 이 $\mu(\theta)$ 를 측정 전위에 결선하는
것이 다음 소절이다.
\end{verifybox}

\begin{figure}[t]
\centering
\begin{tikzpicture}[x=9.6cm,y=6.8cm]
\draw[-{Latex[length=1.6mm]}] (-0.02,-0.76) -- (-0.02,0.14) node[above,font=\scriptsize] {$f(\xi)/RT$ (mixing part)};
\draw[-{Latex[length=1.6mm]}] (-0.02,0) -- (1.06,0) node[below,font=\scriptsize] {$\xi$};
\foreach \xx in {0.5,1.0} {\draw (\xx,0.008) -- (\xx,-0.008) node[below,font=\scriptsize] {\xx};}
% Omega = 0 RT (ideal mixing)
\draw[black!45] plot[smooth] coordinates {(0.004,-0.0261) (0.020,-0.0980) (0.070,-0.2536) (0.118,-0.3625) (0.166,-0.4488) (0.213,-0.5183) (0.261,-0.5742) (0.309,-0.6182) (0.357,-0.6515) (0.404,-0.6748) (0.452,-0.6886) (0.500,-0.6931) (0.548,-0.6886) (0.596,-0.6748) (0.643,-0.6515) (0.691,-0.6182) (0.739,-0.5742) (0.787,-0.5183) (0.834,-0.4488) (0.882,-0.3625) (0.930,-0.2536) (0.980,-0.0980) (0.996,-0.0261)};
% Omega = 1 RT
\draw[black!45] plot[smooth] coordinates {(0.004,-0.0221) (0.020,-0.0784) (0.070,-0.1885) (0.118,-0.2586) (0.166,-0.3106) (0.213,-0.3505) (0.261,-0.3813) (0.309,-0.4047) (0.357,-0.4220) (0.404,-0.4339) (0.452,-0.4409) (0.500,-0.4431) (0.548,-0.4409) (0.596,-0.4339) (0.643,-0.4220) (0.691,-0.4047) (0.739,-0.3813) (0.787,-0.3505) (0.834,-0.3106) (0.882,-0.2586) (0.930,-0.1885) (0.980,-0.0784) (0.996,-0.0221)};
% Omega = 2 RT (threshold: flat bottom)
\draw[thick,black!70] plot[smooth] coordinates {(0.004,-0.0181) (0.020,-0.0588) (0.070,-0.1234) (0.118,-0.1547) (0.166,-0.1725) (0.213,-0.1827) (0.261,-0.1884) (0.309,-0.1913) (0.357,-0.1926) (0.404,-0.1930) (0.452,-0.1931) (0.500,-0.1931) (0.548,-0.1931) (0.596,-0.1930) (0.643,-0.1926) (0.691,-0.1913) (0.739,-0.1884) (0.787,-0.1827) (0.834,-0.1725) (0.882,-0.1547) (0.930,-0.1234) (0.980,-0.0588) (0.996,-0.0181)};
% Omega = 2.5 RT (double well)
\draw[densely dashed,thick] plot[smooth] coordinates {(0.004,-0.0161) (0.020,-0.0490) (0.070,-0.0909) (0.118,-0.1027) (0.166,-0.1034) (0.213,-0.0988) (0.261,-0.0919) (0.309,-0.0845) (0.357,-0.0778) (0.404,-0.0726) (0.452,-0.0693) (0.500,-0.0681) (0.548,-0.0693) (0.596,-0.0726) (0.643,-0.0778) (0.691,-0.0845) (0.739,-0.0919) (0.787,-0.0988) (0.834,-0.1034) (0.882,-0.1027) (0.930,-0.0909) (0.980,-0.0490) (0.996,-0.0161)};
% Omega = 3 RT (double well, humped)
\draw[thick] plot[smooth] coordinates {(0.004,-0.0141) (0.020,-0.0392) (0.070,-0.0583) (0.118,-0.0508) (0.166,-0.0343) (0.213,-0.0149) (0.261,0.0046) (0.309,0.0222) (0.357,0.0369) (0.404,0.0478) (0.452,0.0546) (0.500,0.0569) (0.548,0.0546) (0.596,0.0478) (0.643,0.0369) (0.691,0.0222) (0.739,0.0046) (0.787,-0.0149) (0.834,-0.0343) (0.882,-0.0508) (0.930,-0.0583) (0.980,-0.0392) (0.996,-0.0141)};
% spinodal markers on Omega=3RT: xi_s = 0.2113 / 0.7887, f/RT = -0.0157
\fill (0.2113,-0.0157) circle (0.5pt) node[below left,font=\scriptsize] {$\xi_s^-$};
\fill (0.7887,-0.0157) circle (0.5pt) node[below right,font=\scriptsize] {$\xi_s^+$};
% curve labels at xi = 0.5 (each just off its own curve)
\node[font=\scriptsize] at (0.5,0.095) {$\Omega=3RT$};
\node[font=\scriptsize] at (0.5,-0.115) {$2.5RT$};
\node[font=\scriptsize] at (0.5,-0.235) {$2RT$};
\node[font=\scriptsize] at (0.5,-0.40) {$1RT$};
\node[font=\scriptsize] at (0.5,-0.655) {$\Omega=0$};
\end{tikzpicture}
\caption{정규용액 조성 자유에너지의 섞임 몫 $f(\xi)/RT=\xi\ln\xi+(1-\xi)\ln(1-\xi)+(\Omega/RT)\,\xi(1-\xi)$
(좌표는 식 그대로의 수치 평가). $\Omega\le2RT$ 면 한 웰, $\Omega=2RT$ 에서 바닥이 평탄해지고
(문턱, 식~\eqref{eq:sm-thresh}), $\Omega>2RT$ 면 가운데가 언덕이 된 이중웰 --- $\Omega=3RT$ 곡선의 점이
spinodal $\xi_s^\pm=0.2113/0.7887$ 이다. 닫힌 근은 \S\ref{sec:hys} 가 닫으며, 그 절의 이중웰 그림이 이 가족
중 $\Omega=3RT$ 한 장을 확대한 것이다.}
\label{fig:sm-gxi}
\end{figure}

\begin{figure}[t]
\centering
\begin{tikzpicture}[x=9.6cm,y=0.60cm]
\draw[-{Latex[length=1.6mm]}] (-0.02,0) -- (1.06,0) node[below,font=\scriptsize] {$\theta$};
\draw[-{Latex[length=1.6mm]}] (-0.02,-4.3) -- (-0.02,4.4)
  node[above,font=\scriptsize] {$(\mu-\mu^0)/RT$};
\foreach \xx in {0.25,0.5,0.75} {\draw (\xx,0.07) -- (\xx,-0.07) node[below,font=\scriptsize] {\xx};}
\foreach \yy in {-4,-2,2,4} {\draw (-0.014,\yy) -- (-0.026,\yy) node[left,font=\scriptsize] {\yy};}
% Omega/RT = 0
\draw[thin] plot[smooth] coordinates {(0.02,-3.892) (0.036,-3.288) (0.052,-2.903) (0.068,-2.618) (0.084,-2.389) (0.1,-2.197) (0.13,-1.901) (0.1689,-1.593) (0.2079,-1.338) (0.2468,-1.116) (0.2858,-0.9159) (0.3247,-0.7321) (0.3637,-0.5594) (0.4026,-0.3945) (0.4416,-0.2348) (0.4805,-0.07793) (0.5195,0.07793) (0.5584,0.2348) (0.5974,0.3945) (0.6363,0.5594) (0.6753,0.7321) (0.7142,0.9159) (0.7532,1.116) (0.7921,1.338) (0.8311,1.593) (0.87,1.901) (0.9,2.197) (0.916,2.389) (0.932,2.618) (0.948,2.903) (0.964,3.288) (0.98,3.892)};
% Omega/RT = 2 (threshold)
\draw[densely dashed] plot[smooth] coordinates {(0.02,-1.972) (0.036,-1.432) (0.052,-1.111) (0.068,-0.8898) (0.084,-0.7252) (0.1,-0.5972) (0.13,-0.421) (0.1689,-0.2689) (0.2079,-0.1692) (0.2468,-0.1029) (0.2858,-0.05908) (0.3247,-0.03103) (0.3637,-0.01415) (0.4026,-0.005038) (0.4416,-0.001072) (0.4805,-3.942e-05) (0.5195,3.942e-05) (0.5584,0.001072) (0.5974,0.005038) (0.6363,0.01415) (0.6753,0.03103) (0.7142,0.05908) (0.7532,0.1029) (0.7921,0.1692) (0.8311,0.2689) (0.87,0.421) (0.9,0.5972) (0.916,0.7252) (0.932,0.8898) (0.948,1.111) (0.964,1.432) (0.98,1.972)};
% Omega/RT = 4 (van der Waals loop)
\draw[thick] plot[smooth] coordinates {(0.02,-0.05182) (0.036,0.4244) (0.052,0.6809) (0.068,0.8382) (0.084,0.9388) (0.1,1.003) (0.13,1.059) (0.1689,1.055) (0.2079,0.9992) (0.2468,0.9097) (0.2858,0.7978) (0.3247,0.67) (0.3637,0.5311) (0.4026,0.3844) (0.4416,0.2326) (0.4805,0.07786) (0.5195,-0.07786) (0.5584,-0.2326) (0.5974,-0.3844) (0.6363,-0.5311) (0.6753,-0.67) (0.7142,-0.7978) (0.7532,-0.9097) (0.7921,-0.9992) (0.8311,-1.055) (0.87,-1.059) (0.9,-1.003) (0.916,-0.9388) (0.932,-0.8382) (0.948,-0.6809) (0.964,-0.4244) (0.98,0.05182)};
% extrema (= spinodal) on Omega/RT=4
\fill (0.1464,1.0657) circle (0.5pt) node[above,font=\scriptsize] {$\theta_s^-$};
\fill (0.8536,-1.0657) circle (0.5pt) node[below,font=\scriptsize] {$\theta_s^+$};
% curve labels
\node[font=\scriptsize] at (0.115,-3.05) {$\Omega=0$};
\node[font=\scriptsize,right] at (0.99,1.97) {$\Omega=2RT$};
\node[font=\scriptsize] at (0.10,1.75) {$\Omega=4RT$};
\end{tikzpicture}
\caption{격자기체 화학퍼텐셜 $(\mu(\theta)-\mu^0)/RT=\ln[\theta/(1-\theta)]+(\Omega/RT)(1-2\theta)$
의 개형(식~\eqref{eq:mu}; 좌표는 식 그대로의 수치 평가). $\Omega=0$ 단조(로그 몫이 양끝을
잠금), $\Omega=2RT$ 에서 중심 기울기 $0$(문턱, 식~\eqref{eq:sm-thresh}), $\Omega=4RT$ 는 비단조 --- 한
$\mu$ 에 세 $\theta$ 가 대응하는 van der Waals loop 이고 극값(점)이 spinodal $\theta_s^\pm=0.146/0.854$,
값 $\pm1.066$. \S\ref{sec:hys} 과주행 그림의 $V_\eq(\xi)$ 곡선은 이 곡선의 $\theta=1-\xi$·부호 반전
거울이라 같은 $\pm1.066$ 이 나타난다.}
\label{fig:sm-mu}
\end{figure}

\subsection{전기화학 연결 --- $\mu_\mathrm{Li}=\mu^0-sF(V-U)$ 와 평형 점유 logistic}\label{sec:sm-electro}
\textbf{(a) 출발 --- 전기화학 퍼텐셜과 삽입 반응.} 하전 종은 화학퍼텐셜에 정전 에너지가 더해진
\emph{전기화학 퍼텐셜}(electrochemical potential)로 평형을 판정한다:
\begin{equation}
\tilde\mu\;\equiv\;\mu+zF\phi
\label{eq:sm-emu}
\end{equation}
($z$ = 전하수, $\phi$ = 그 종이 있는 상의 정전 퍼텐셜). 삽입 반응
$\mathrm{Li}^++e^-+[\,]\rightleftharpoons\mathrm{Li}_{(\mathrm{host})}$ 의 평형은 양변 $\tilde\mu$ 합의 일치이고
(\S\ref{sec:center} 가 결과 사슬에서 같은 균형으로 진입한다), host 안의 Li 는 중성이라
$\tilde\mu_\mathrm{Li}=\mu_\mathrm{Li}$ 다:
\begin{equation}
\big(\mu_{\mathrm{Li}^+}+F\phi_S\big)+\big(\mu_{e^-}-F\phi_M\big)=\mu_\mathrm{Li}(\theta)
\label{eq:sm-workbal}
\end{equation}
($\phi_S$ = 전해질, $\phi_M$ = 작동 전극의 정전 퍼텐셜; $z_{\mathrm{Li}^+}{=}+1$, $z_{e^-}{=}-1$).

\textbf{(b) 연산 --- 기준전극으로 측정 전위를 결선.} 같은 전해질에 담긴 기준 Li 금속 전극도 자기 평형
$\mathrm{Li}^++e^-\rightleftharpoons\mathrm{Li}_{(\mathrm{metal})}$ 을 이룬다:
\begin{equation}
\big(\mu_{\mathrm{Li}^+}+F\phi_S\big)+\big(\mu_{e^-}-F\phi_\mathrm{ref}\big)=\mu_\mathrm{Li}^\mathrm{metal}.
\label{eq:sm-refbal}
\end{equation}
\eqref{eq:sm-workbal} 에서 \eqref{eq:sm-refbal} 을 빼면 전해질 몫($\mu_{\mathrm{Li}^+}+F\phi_S$)이
상쇄된다. 전자의 화학 몫도 상쇄되는데, 여기에는 한 단계가 숨어 있다 --- 두 전극의 전자 화학퍼텐셜은
엄밀히는 재질에 따라 다르며, 그 차이(접촉 전위차)는 두 전극을 같은 재질의 도선 단자 사이에서 읽는 측정
전위의 조작적 정의에 흡수되어 별도 항으로 남지 않는다(전기화학의 표준 처리). 그 결과 측정 전위
$V\equiv\phi_M-\phi_\mathrm{ref}$(vs Li/Li$^+$)만 남는다:
\begin{equation}
\mu_\mathrm{Li}(\theta_\eq)-\mu_\mathrm{Li}^\mathrm{metal}=-F\,V.
\label{eq:sm-muV}
\end{equation}
전위를 올리면 host 안 Li 의 화학퍼텐셜이 정확히 $-FV$ 만큼 낮아진다 --- ``$V$ 를 올리면 탈리튬화''가 이 한 줄이다.

\textbf{(c) 중간식 --- 상수 덩이를 $U$ 로 묶기.} 식~\eqref{eq:sm-muV} 의 상수 $\mu_\mathrm{Li}^\mathrm{metal}$
를 전이 중심 $U\equiv(\mu_\mathrm{Li}^\mathrm{metal}-\mu^0)/F$ 로 재포장하면($s=+1$ --- \S\ref{sec:notation}
표의 유도 전용 고정 부호, 방향 부호 $\sigma_d$ 와 별개)
\begin{equation}
\mu_\mathrm{Li}(\theta_\eq)=\mu^0-sF\,(V-U),
\qquad
\Delta G_j\equiv\mu^0-\mu_\mathrm{Li}^\mathrm{metal}=-sF\,U_j
\label{eq:sm-eqcond}
\end{equation}
--- \S\ref{sec:center}(N2)가 결과 사슬에서 다시 세우는 평형 조건과 글자까지 같은 식이다($U_j>0
\Leftrightarrow\Delta G_j<0$, 곧 그 전이가 자발일 전위 창이 양수라는 부호 읽기도 동일). Chapter 2 의 식
(eq:muV)도 같은 관계다.

\textbf{(d) 박스 --- 평형 점유 logistic.} $\Omega_j=0$ 이면 \eqref{eq:sm-muideal} 를 \eqref{eq:sm-eqcond} 에
넣어 $RT\ln[\theta_\eq/(1-\theta_\eq)]=-sF(V-U)$, 로그를 풀면 닫힌 꼴이 나온다:
\begin{equation}
\boxed{\;\theta_\eq(V,T)=\frac{1}{1+e^{+sF(V-U)/RT}},\qquad
\xi_\eq(V,T)=1-\theta_\eq=\frac{1}{1+e^{-sF(V-U)/RT}},\qquad
w\equiv\frac{RT}{F}\;.}
\label{eq:sm-logistic}
\end{equation}
자리당 언어와의 일치 검산 --- 식~\eqref{eq:fermifn} 의 지수는 $\beta\,\Delta\mu$ 이고, 몰로 올리면
$(\mu^0-\mu_\mathrm{Li})/RT$, 여기에 \eqref{eq:sm-eqcond} 를 넣으면 $+sF(V-U)/RT$ --- 정확히
\eqref{eq:sm-logistic} 의 $\theta_\eq$ 지수다(자리당 쌍 $(k_BT,\,e)$ 가 몰당 $(RT,\,F)$ 로 한꺼번에 환산).
점유 $\theta_\eq$ 는 $V$ 에 감소하고 탈리튬화 진행률 $\xi_\eq$ 는 증가하는 logistic 이며,
그림~\ref{fig:sm-occ} 가 온도별 개형과 두 곡선의 여집합 교차를 보인다. 여기서 세 가지가 본론으로 이어진다.
\begin{enumerate}[label=(\roman*),leftmargin=2.2em,itemsep=1pt]
\item \textbf{일반화는 \S\ref{sec:width} 소관.} 고정 부호 $s$ 자리에 방향 부호 $\sigma_d$, 중심에 분기 중심
$U_j^{\,d}$(\S\ref{sec:hys}), 폭에 다중도 $w_j=n_jRT/F$ 를 넣은 것이 모델의 평형 진행률 $\xi_\eq$ 다 --- 단 지수의
$\sigma_d$ 는 $\mathrm{logistic}[-x]=1-\mathrm{logistic}[x]$ 에 의한 $\xi\leftrightarrow1-\xi$ 여집합
relabel 일 뿐 종 $\xi(1-\xi)$ 은 불변이라, 물리 작용처(\S\ref{sec:notation} 의 분극$\cdot$분기$\cdot$꼬리)로는
세지 않는다.
\item \textbf{동역학 경로와의 합류.} 같은 logistic 이 Eyring 속도식의 정지점(detailed balance,
\S\ref{sec:width})에서도 나온다 --- 평형비 $\xi_\eq/(1-\xi_\eq)=e^{sF(V-U)/RT}$ 가 두 경로에서 같으니,
통계(이 절)와 동역학(\S\ref{sec:width})은 한 분포의 두 유도다(\S\ref{sec:dist} 의 명제를 Part 0 이 평형
쪽에서 선결).
\item \textbf{$\Omega_j\ne0$ 이면 닫힌 logistic 이 아니다.} \eqref{eq:sm-eqcond} 에 상호작용까지 든
식~\eqref{eq:mu} 를 넣고 $\xi$ 좌표의 우함수 대칭(\S\ref{sec:sm-mf})을 쓰면
$RT\ln[\xi/(1-\xi)]+\Omega_j(1-2\xi)=sF(V-U)$ 의 암시적 등온선이 되고(\S\ref{sec:hys} 의 $V_\eq(\xi)$),
$\Omega_j>2RT$(식~\eqref{eq:sm-thresh})면 비단조다. 곧 닫힌 logistic 은 $\Omega_j=0$ 에서만 엄밀하며,
$\Omega_j\ne0$ 전이에 같은 서식을 계속 쓰는 것의 지위(평형 예측 vs 현상학적 폭)는 \S\ref{sec:width} 의
폭 이중지위가 가른다.
\end{enumerate}

\begin{signbox}
\textbf{부호(Part 0).} $s=+1$: $V\!\uparrow\Rightarrow\mu_\mathrm{Li}\!\downarrow$(식~\eqref{eq:sm-eqcond})
$\Rightarrow\theta_\eq\!\downarrow,\ \xi_\eq\!\uparrow$ --- 전위를 올리면 탈리튬화가 진행한다(본론 규약 일치).
$V=U\Rightarrow\theta_\eq=\xi_\eq=\tfrac12$, $T\to0$ 계단($w\to0$)·$T\!\uparrow$ 폭 비례($w=RT/F$).
방향 부호 $\sigma_d$ 의 작용처(분극$\cdot$분기$\cdot$꼬리)는 \S\ref{sec:notation}·\S\ref{sec:signcheck} 소관 --- Part 0 은 $s=+1$ 고정.
\end{signbox}

\begin{figure}[t]
\centering
\resizebox{\textwidth}{!}{%
\begin{tikzpicture}
\begin{scope}[x=0.50cm,y=2.4cm]
% ---- (a) dimensionless single-site occupancy ----
\draw[-{Latex[length=1.6mm]}] (-6.6,0) -- (7.0,0)
  node[below,font=\scriptsize] {$(\varepsilon-\mu)/k_BT_0$};
\draw[-{Latex[length=1.6mm]}] (-6.6,0) -- (-6.6,1.18) node[above,font=\scriptsize] {$\theta$};
\foreach \xx in {-4,-2,2,4} {\draw (\xx,0.013) -- (\xx,-0.013) node[below,font=\scriptsize] {\xx};}
\draw (-6.54,1.0) -- (-6.66,1.0) node[left,font=\scriptsize] {1};
\draw (-6.54,0.5) -- (-6.66,0.5) node[left,font=\scriptsize] {0.5};
\draw[densely dotted,gray] (-6.6,0.5) -- (6.6,0.5);
% T = T0/2
\draw[thick,densely dotted] plot[smooth] coordinates {(-6,1) (-5,1) (-4,0.9997) (-3.5,0.9991) (-3,0.9975) (-2.5,0.9933) (-2,0.982) (-1.5,0.9526) (-1,0.8808) (-0.5,0.7311) (0,0.5) (0.5,0.2689) (1,0.1192) (1.5,0.04743) (2,0.01799) (2.5,0.006693) (3,0.002473) (3.5,0.0009111) (4,0.0003354) (5,4.54e-05) (6,6.144e-06)};
% T = T0
\draw[thick] plot[smooth] coordinates {(-6,0.9975) (-5.5,0.9959) (-5,0.9933) (-4.5,0.989) (-4,0.982) (-3.5,0.9707) (-3,0.9526) (-2.5,0.9241) (-2,0.8808) (-1.5,0.8176) (-1,0.7311) (-0.5,0.6225) (0,0.5) (0.5,0.3775) (1,0.2689) (1.5,0.1824) (2,0.1192) (2.5,0.07586) (3,0.04743) (3.5,0.02931) (4,0.01799) (4.5,0.01099) (5,0.006693) (5.5,0.00407) (6,0.002473)};
% T = 2T0
\draw[thick,densely dashed] plot[smooth] coordinates {(-6,0.9526) (-5.5,0.9399) (-5,0.9241) (-4.5,0.9047) (-4,0.8808) (-3.5,0.852) (-3,0.8176) (-2.5,0.7773) (-2,0.7311) (-1.5,0.6792) (-1,0.6225) (-0.5,0.5622) (0,0.5) (0.5,0.4378) (1,0.3775) (1.5,0.3208) (2,0.2689) (2.5,0.2227) (3,0.1824) (3.5,0.148) (4,0.1192) (4.5,0.09535) (5,0.07586) (5.5,0.06009) (6,0.04743)};
\node[font=\scriptsize,align=left] at (3.6,0.82)
  {dotted: $T=T_0/2$\\ solid: $T=T_0$\\ dashed: $T=2T_0$};
\node[font=\scriptsize] at (0,-0.30) {(a)};
\end{scope}
\begin{scope}[xshift=8.9cm,x=33cm,y=2.4cm]
% ---- (b) electrochemical dress: xi_eq(V), theta_eq(V) ----
\draw[-{Latex[length=1.6mm]}] (-0.006,0) -- (0.215,0) node[below,font=\scriptsize] {$V$ [V]};
\draw[-{Latex[length=1.6mm]}] (-0.006,0) -- (-0.006,1.18)
  node[above,font=\scriptsize] {$\xi_\eq,\ \theta_\eq$};
\foreach \xx in {0.05,0.1,0.15,0.2}
  {\draw (\xx,0.013) -- (\xx,-0.013) node[below,font=\scriptsize] {\xx};}
\draw[densely dotted] (0.085,0) -- (0.085,1.0);
\node[below,font=\scriptsize] at (0.085,-0.10) {$U$};
% xi_eq, T=268.15 K
\draw[thick,densely dotted] plot[smooth] coordinates {(0,0.02464) (0.008,0.03448) (0.016,0.04806) (0.024,0.06662) (0.032,0.09165) (0.04,0.1248) (0.048,0.1678) (0.056,0.2218) (0.064,0.2872) (0.072,0.3629) (0.08,0.4461) (0.088,0.5324) (0.096,0.6168) (0.104,0.6947) (0.112,0.7629) (0.12,0.8198) (0.128,0.8654) (0.136,0.9009) (0.144,0.9278) (0.152,0.9478) (0.16,0.9625) (0.168,0.9732) (0.176,0.9809) (0.184,0.9864) (0.192,0.9903) (0.2,0.9932)};
% xi_eq, T=298.15 K
\draw[thick] plot[smooth] coordinates {(0,0.03529) (0.008,0.04756) (0.016,0.06383) (0.024,0.08516) (0.032,0.1128) (0.04,0.1479) (0.048,0.1915) (0.056,0.2444) (0.064,0.3063) (0.072,0.3761) (0.08,0.4515) (0.088,0.5292) (0.096,0.6054) (0.104,0.6769) (0.112,0.7409) (0.12,0.7961) (0.128,0.8421) (0.136,0.8792) (0.144,0.9086) (0.152,0.9314) (0.16,0.9488) (0.168,0.962) (0.176,0.9719) (0.184,0.9792) (0.192,0.9847) (0.2,0.9887)};
% xi_eq, T=328.15 K
\draw[thick,densely dashed] plot[smooth] coordinates {(0,0.04716) (0.008,0.06163) (0.016,0.08017) (0.024,0.1037) (0.032,0.133) (0.04,0.1692) (0.048,0.2127) (0.056,0.2639) (0.064,0.3224) (0.072,0.3871) (0.08,0.4559) (0.088,0.5265) (0.096,0.596) (0.104,0.6619) (0.112,0.7221) (0.12,0.7752) (0.128,0.8206) (0.136,0.8586) (0.144,0.8896) (0.152,0.9145) (0.16,0.9342) (0.168,0.9496) (0.176,0.9615) (0.184,0.9707) (0.192,0.9778) (0.2,0.9832)};
% theta_eq = 1 - xi_eq, T=298.15 K
\draw[gray,thick] plot[smooth] coordinates {(0,0.9647) (0.008,0.9524) (0.016,0.9362) (0.024,0.9148) (0.032,0.8872) (0.04,0.8521) (0.048,0.8085) (0.056,0.7556) (0.064,0.6937) (0.072,0.6239) (0.08,0.5485) (0.088,0.4708) (0.096,0.3946) (0.104,0.3231) (0.112,0.2591) (0.12,0.2039) (0.128,0.1579) (0.136,0.1208) (0.144,0.09142) (0.152,0.06864) (0.16,0.05122) (0.168,0.03803) (0.176,0.02814) (0.184,0.02077) (0.192,0.0153) (0.2,0.01125)};
\node[font=\scriptsize] at (0.163,0.72) {$\xi_\eq$};
\node[gray,font=\scriptsize] at (0.033,0.70) {$\theta_\eq=1-\xi_\eq$};
\node[font=\scriptsize,align=left] at (0.175,0.30)
  {dotted: 268 K\\ solid: 298 K\\ dashed: 328 K};
\node[font=\scriptsize] at (0.1,-0.30) {(b)};
\end{scope}
\end{tikzpicture}%
}
\caption{단일 자리 점유의 전기화학적 재매개변수화(좌표는 식 그대로의 수치 평가). (a) 식~\eqref{eq:fermifn} 의
$\theta$ --- $\varepsilon=\mu$ 에서 $\tfrac12$, 낮은 $T$ 일수록 계단($\beta\to\infty$ --- $T\to0$ 극한이
계단 함수다), 높은 $T$ 일수록 완만. (b) 식~\eqref{eq:sm-logistic} 의 $\xi_\eq(V)$ ($U=0.085$ V, stage
$2\!\to\!1$ 초기값) --- 폭 $w=RT/F$ 가 $23.1/25.7/28.3$ mV($268/298/328$ K)로 $T$ 에 비례해 열리고, 회색
$\theta_\eq=1-\xi_\eq$ 와 중심 $V=U$ 에서 교차한다($\tfrac12$).}
\label{fig:sm-occ}
\end{figure}

\subsection{거시 열역학으로 --- $G\equiv H-TS$, $\Delta G_j=-sFU_j$, Nernst 식}\label{sec:sm-macro}
\textbf{(a) 출발 --- 두 언어의 같은 $\mu$.} 본론 \S\ref{sec:center} 는 거시 정의 --- Gibbs 자유에너지
$G\equiv H-TS$ 와 $\mu\equiv\partial G/\partial n|_{T,P}$ --- 에서 출발한다. 통계역학 쪽 정의는
\S\ref{sec:sm-ensemble} 의 $\mu=\partial F/\partial N|_{T,V}$ 였다. 둘의 다리는 $G=F+PV$ 다 --- 응축상
격자에 이온 1몰을 넣을 때의 $P\Delta v$ 몫은 $RT$·전기 일($FV$) 스케일에 비해 무시되므로
\begin{equation}
\mu=\frac{\partial G}{\partial n}\Big|_{T,P}\;\simeq\;\frac{\partial F}{\partial n}\Big|_{T,V}
\qquad(\text{삽입 고체: }P\Delta v\ll RT,\;F\,V)
\label{eq:sm-mubridge}
\end{equation}
--- \S\ref{sec:sm-lattice}--\S\ref{sec:sm-electro} 에서 세운 $\mu(\theta)$ 를 본론의 거시 $\mu$ 자리에
그대로 꽂아도 되는 근거가 이 한 줄이다.

\textbf{(b) 연산 --- 반응 자유에너지의 $H,S$ 분해.} \eqref{eq:sm-eqcond} 의 비배치 몫
$\Delta G_j=\mu^0-\mu_\mathrm{Li}^\mathrm{metal}$ 은 온도의 함수이고, $G\equiv H-TS$ 를 반응 차분에 적용하면
$\Delta G_j(T)=\Delta H_{\rxn,j}-T\Delta S_{\rxn,j}$ 다.

\textbf{(c) 중간식 --- 평형 중심의 온도 환산.}
이를 $\Delta G_j=-sFU_j$ 와 등치하면 $\Delta H_{\rxn,j}-T\Delta S_{\rxn,j}=-FU_j(T)$($s{=}1$) --- 그 이항이
\S\ref{sec:center} 의 결과 박스($U_j(T)$·$\partial U_j/\partial T=\Delta S_{\rxn,j}/F$)이며, 박스는 그 절의
것을 그대로 쓰고 여기 다시 적지 않는다.

\textbf{(d) 박스 --- Nernst 식으로 마감.} 마지막으로 \eqref{eq:sm-logistic} 의 $\xi_\eq$ 를 전위에 대해
뒤집으면($\xi/(1-\xi)=e^{sF(V-U)/RT}$ 의 로그) 이상 극한($\Omega_j=0$)의 평형 전위 곡선이 닫힌다:
\begin{equation}
\boxed{\;
\Delta G_j=-sF\,U_j,
\qquad
V_\eq(\xi)=U_j+\frac{RT}{sF}\,\ln\frac{\xi}{1-\xi}
\qquad(\Omega_j=0,\ s=+1)\;.}
\label{eq:sm-nernst}
\end{equation}
로그 몫의 정체는 \eqref{eq:sm-Smix} 섞임 엔트로피의 부분몰 몫이다 --- Nernst 식의 $\ln[\xi/(1-\xi)]$ 는
곡선맞춤이 아니라 배치 셈법에서 온다(Chapter 2 의 식 (eq:Vxi)가 같은 결론을 발열 관점에서 재사용한다).
\begin{verifybox}
극한·부호 검산 --- $\xi=\tfrac12$ 에서 $V=U_j$(중심), $T\!\uparrow$ 이면 로그 몫의 기울기(폭 $w=RT/F$)가
비례해 커지고, $\Delta S_{\rxn,j}>0$ 이면 중심이 온도에 오른다. 상호작용 몫 $\Omega_j(1-2\xi)/sF$ 를 더하면
\S\ref{sec:hys} 의 $V_\eq(\xi)$ 로 확장되고, $\Omega_j>2RT$ 의 비단조가 히스테리시스로 이어진다.
\end{verifybox}

\begin{keybox}
\textbf{Part 0 사다리.} 등확률~\eqref{eq:sm-S} $\to$ Gibbs 인자·$\Xi$~\eqref{eq:sm-gc} $\to$ 단일 자리
$\theta$~\eqref{eq:fermifn}·요동~\eqref{eq:sm-flucres} $\to$ lattice gas·섞임 엔트로피~\eqref{eq:sm-Smix}
$\to$ 평균장 $\Omega$~\eqref{eq:mu}·\eqref{eq:gxi}·문턱~\eqref{eq:sm-thresh} $\to$ 전위
손잡이~\eqref{eq:sm-eqcond} $\to$ logistic~\eqref{eq:sm-logistic}·Nernst~\eqref{eq:sm-nernst}. 본론이 결과로
쓰는 평형식의 전 유도가 이 사다리에 있다 --- 이후 절은 같은 사다리를 결과 사슬(N1--N9) 순서로
다시 밟는다.
\end{keybox}

% 자산: [A-047] [A-048] [A-049] [A-050] [A-051] [A-052] [A-053] [A-054] [A-055] [A-056] [A-057] [A-058] [A-059] [A-060] [A-061] [A-062] [A-063] [A-064] [A-065] [A-066] [A-067] [A-068] [A-069]
